\documentclass{article}
\usepackage[margin=1in]{geometry}
\usepackage[autocompile]{gregoriotex}

\newcommand{\lines}[1]{
  \begin{minipage}[b]{1.5cm}
    \grechangedim{parskip}{0pt}{fixed}%
    \setlength{\baselineskip}{0pt}%
    \setlength{\lineskip}{0pt}%
    \hrule
    \setbox0=\vbox{#1}%
    \copy0%
    \hrule
    \vspace{-\dp0}%
    \hrule
  \end{minipage}%
}

% same as \gabcsnippet, but with nabc lines turned on (see GregorioRef.tex)
\gresetnabcfont[1]{gregall}{16}
\long\def\nabcsnippet#1#2{%
  \directlua{gregoriotex.direct_gabc("\luaescapestring{\unexpanded\expandafter{#2}}",
                                     "nabc-lines:#1;", false)}%
}%

\setlength{\parindent}{0pt}

\begin{document}
\gresetinitiallines{0}
\gresetclef{invisible}

% Above-lines nabc has its own threshold (additionaltopspacenabcthreshold),
% same idea as additionaltopspacealtthreshold (see vertical-spacing.tex) but
% for nabc instead of alt text.
%
% \gre@vepisemaorrare clamps a rare sign's height (accentus, circulus,
% semicirculus, musica ficta) at \gre@pitch@raresign (= \gre@pitch@abovestaff),
% a pitch fixed by the staff-lines count -- so a rare sign never renders
% below that floor, no matter how low the note it's attached to. That floor
% sits exactly one pitch step above "adjust_top", the reference
% gregoriotex.lua compares a line's content against to decide whether
% above-lines nabc needs to grow. A threshold higher than 0 can never see a
% one-step difference, so lowering it to 0 (the default) makes above-lines
% nabc react to a rare sign at exactly the pitches it always renders at,
% regardless of how low the underlying note actually is.
%
% That reactivity alone doesn't fully close the gap, though: the accentus
% glyph's tip extends a bit higher than the pitch-step reservation assumes,
% so above-lines nabc can still end up close enough to touch it at the top
% of the range (see the grid below -- a and c have plenty of room, but the
% gap narrows steadily up to m). This was deliberately left as a case-by-case
% manual adjustment (\grechangedim{abovelinesnabcraise}{...}{scalable}, see
% below) rather than an automatic default-value increase, since bumping
% abovelinesnabcraise unconditionally would add the same extra headroom to
% every line, including ones nowhere near this collision.
{
\verb|\grechangecount{additionaltopspacenabcthreshold}{0} (default)|
\grechangecount{additionaltopspacenabcthreshold}{0}

\begin{center}
\begin{tabular}{ccccccc}
  \lines{\nabcsnippet{1}{a(ar1|vi)}} &
  \lines{\nabcsnippet{1}{c(cr1|vi)}} &
  \lines{\nabcsnippet{1}{e(er1|vi)}} &
  \lines{\nabcsnippet{1}{g(gr1|vi)}} &
  \lines{\nabcsnippet{1}{i(ir1|vi)}} &
  \lines{\nabcsnippet{1}{k(kr1|vi)}} &
  \lines{\nabcsnippet{1}{m(mr1|vi)}}
\end{tabular}
\end{center}
}

% Same collision, forcing additionaltopspacenabcthreshold back to 4, its
% value before #1766 was fixed. Every one of these notes carries the same
% rare sign as the row above; the only thing that changed is the threshold.
% This reproduces the pre-fix bug in its most severe form: above-lines nabc
% never grows to make room for the rare sign at all, no matter how close the
% sign's rendered height gets to it, since 4 is more than the one-step
% difference the floor can ever produce over adjust_top. This is also the
% way to opt back into the old fixed-distance behavior, if for some reason
% it is preferred.
{
\verb|\grechangecount{additionaltopspacenabcthreshold}{4} (opts out)|
\grechangecount{additionaltopspacenabcthreshold}{4}

\begin{center}
\begin{tabular}{ccccccc}
  \lines{\nabcsnippet{1}{a(ar1|vi)}} &
  \lines{\nabcsnippet{1}{c(cr1|vi)}} &
  \lines{\nabcsnippet{1}{e(er1|vi)}} &
  \lines{\nabcsnippet{1}{g(gr1|vi)}} &
  \lines{\nabcsnippet{1}{i(ir1|vi)}} &
  \lines{\nabcsnippet{1}{k(kr1|vi)}} &
  \lines{\nabcsnippet{1}{m(mr1|vi)}}
\end{tabular}
\end{center}
}

% The manual workaround for the remaining collision at the top of the
% default grid above: raise abovelinesnabcraise for the affected score.
% 0.25cm is the value this branch shipped as the new default at one point,
% before that unconditional change was reverted in favor of leaving it to
% be tuned per score, since it isn't needed on lines without a rare sign or
% very high note.
{
\verb|\grechangedim{abovelinesnabcraise}{0.25cm}{scalable} (manual workaround)|
\grechangedim{abovelinesnabcraise}{0.25cm}{scalable}

\begin{center}
\begin{tabular}{ccccccc}
  \lines{\nabcsnippet{1}{a(ar1|vi)}} &
  \lines{\nabcsnippet{1}{c(cr1|vi)}} &
  \lines{\nabcsnippet{1}{e(er1|vi)}} &
  \lines{\nabcsnippet{1}{g(gr1|vi)}} &
  \lines{\nabcsnippet{1}{i(ir1|vi)}} &
  \lines{\nabcsnippet{1}{k(kr1|vi)}} &
  \lines{\nabcsnippet{1}{m(mr1|vi)}}
\end{tabular}
\end{center}
\grechangedim{abovelinesnabcraise}{0.17351cm}{scalable}
}

% The same floor also catches a plain note with no rare sign at all, once
% its own pitch reaches the same height (h and i sit below it, j one step
% below, k is the floor itself, l and m one and two steps above it) --
% confirming this was never a rare-sign-only bug, just one that #1766
% happened to be filed against. As above, l and m still end up close enough
% to touch above-lines nabc at the default abovelinesnabcraise; the same
% manual workaround applies regardless of whether the tall content is a
% rare sign or a plain note.
{
\verb|\grechangecount{additionaltopspacenabcthreshold}{0} (default, no rare sign)|
\grechangecount{additionaltopspacenabcthreshold}{0}

\begin{center}
\begin{tabular}{cccccc}
  \lines{\nabcsnippet{1}{h(h|vi)}} &
  \lines{\nabcsnippet{1}{i(i|vi)}} &
  \lines{\nabcsnippet{1}{j(j|vi)}} &
  \lines{\nabcsnippet{1}{k(k|vi)}} &
  \lines{\nabcsnippet{1}{l(l|vi)}} &
  \lines{\nabcsnippet{1}{m(m|vi)}}
\end{tabular}
\end{center}
}

% Real-world sanity check: a full antiphon (nabc voice 1) whose melody
% reaches up to j/k ("Dó-mi-num"), at the default
% additionaltopspacenabcthreshold/abovelinesnabcraise, to confirm the fix
% also holds outside the synthetic single-note grids above.
\grechangecount{additionaltopspacenabcthreshold}{0}
\grechangedim{abovelinesnabcraise}{0.17351cm}{scalable}

\section*{Real-world example: Adorate Dominum (mode 8)}

\gregorioscore[f]{adorate-dominum}

% Same antiphon, but with the "mi" neume (over the highest note in the
% line, pitch k) forced down via an explicit nabc pitch code (viha instead
% of vi): ha is the lowest relative pitch on the h-code scale, so this
% pulls that one glyph's ink toward the staff/note independently of the
% line-wide additionaltopspacenabcthreshold/abovelinesnabcraise elevation,
% which only reacts to the underlying square-note pitch and has no
% visibility into this kind of per-glyph offset. Checks whether a lowered
% glyph can still collide with the note it sits above even when the
% line's uniform nabc_raise looks sufficient for the unmodified neumes.
\section*{Same, with mi's neume forced down (viha instead of vi)}

\gregorioscore[f]{adorate-dominum-ha}

\end{document}
